\section{Dream Content: Sensory Modalities and Thematic Elements}\label{sec:content}

Beyond the \textit{form} of dream states, the \textit{content} of human dreams provides crucial data. Large-scale diary and questionnaire studies reveal consistent patterns in both the sensory modalities and thematic elements of our dreams \citep{nielsen2003typical}.

\subsubsection{The Sensory Landscape of Dreams}
Not all senses are represented equally. Dream reports show a clear hierarchy of sensory engagement:
\begin{itemize}
    \item \textbf{Vision and Sound Dominate:} Most dreams are audiovisual experiences. Vision is near-ubiquitous, and sound is also highly common, reported in over 50\% of dreams \citep{zadra2014thematic}.
    \item \textbf{Touch and Movement:} Proprioception (sense of body position) and vestibular sensations (balance, flying, falling) are also frequently engaged, forming the core of many ``typical'' dream themes \citep{nielsen2003typical}.
    \item \textbf{Smell and Taste are Rare:} In stark contrast, olfactory (smell) and gustatory (taste) sensations are exceptionally rare in spontaneous dream reports, often appearing in less than 1\% of dreams \citep{zadra1998prevalence}.
\end{itemize}
This sensory imbalance suggests that dreaming preferentially simulates audiovisual and kinesthetic systems, while largely ignoring the chemical senses.

\subsubsection{The Thematic Universe: Ancient vs. Modern}
Dream themes also fall into distinct categories that reflect both evolved predispositions and modern personal concerns (the ``continuity hypothesis'') \citep{schredl2010continuity}.
\begin{itemize}
    \item \textbf{Universal \& Ancient Themes:} Many of the most common ``typical'' dreams appear to be evolutionarily ancient. Themes like \textbf{being chased}, \textbf{falling}, \textbf{flying}, and \textbf{teeth falling out} are reported across cultures and are often interpreted as rehearsing ancestral survival threats or biological challenges (e.g., via Threat Simulation Theory) \citep{revonsuo2000threat}.
    \item \textbf{Modern \& Continuity Themes:} Other common themes are clearly drawn from daily life in an industrial or post-industrial society. These include \textbf{failing an exam}, \textbf{being late} for an event, \textbf{losing control of a vehicle}, or being unable to find a toilet \citep{nielsen2003typical, schredl2010continuity}. Interestingly, while the \textit{anxiety} of lateness is common, literal clocks are rarely seen \citep{schredl2021clocks}. These dreams reflect the ``continuity'' of our waking anxieties, social-evaluative pressures, and daily logistics.
\end{itemize}